\documentclass[12pt,a4paper]{article}
\usepackage{amsmath,amssymb}
\title{Heat Flow in a Thin Rod}
\author{Leopold Ashgrove}
\date{February 27, 2018}
\begin{document}
\maketitle

\begin{abstract}
We derive the one-dimensional heat equation from conservation of energy and
describe the behaviour of its simplest solutions.
\end{abstract}

\section{Derivation}
Consider a thin rod of length $L$ with temperature $u(x,t)$. Energy
conservation together with Fourier's law gives
\begin{equation}\label{eq:heat}
\frac{\partial u}{\partial t} = \kappa \frac{\partial^2 u}{\partial x^2},
\end{equation}
where $\kappa > 0$ is the thermal diffusivity. The total heat content is
$Q(t) = c\rho \int_0^L u(x,t)\,dx$.

\subsection{Separated Solutions}
Seeking solutions $u = X(x)T(t)$ leads to
\begin{align*}
X'' + \lambda X &= 0, \\
T' + \kappa\lambda T &= 0, \\
X(0) = X(L) &= 0.
\end{align*}
The admissible eigenvalues are $\lambda_n = (n\pi/L)^2$ for $n = 1, 2, \dots$

\section{General Solution}
Superposing the separated solutions, we obtain
\[
u(x,t) = \sum_{n=1}^{\infty} b_n \sin\frac{n\pi x}{L}\, e^{-\kappa n^2\pi^2 t/L^2},
\]
with coefficients determined by the initial data:
\begin{equation}
b_n = \frac{2}{L}\int_0^L u(x,0)\sin\frac{n\pi x}{L}\,dx.
\end{equation}
Equation~\eqref{eq:heat} is thus solved term by term.

\section{Behaviour}
\begin{itemize}
\item[$\bullet$] High frequencies decay fastest.
\item[(a)] The temperature approaches zero everywhere as $t \to \infty$.
\item[(b)] The dominant late-time profile is the first mode,
  $b_1 \sin(\pi x/L)$.
\end{itemize}
The characteristic relaxation time is $\tau = L^2/(\kappa\pi^2)$, so
doubling the rod's length quadruples the time to cool.
\end{document}
